\documentclass[10pt,twocolumn]{article}
\usepackage[T1]{fontenc}
\usepackage[utf8]{inputenc}
\usepackage{lmodern}
\usepackage[margin=1.8cm,top=2.4cm,bottom=2cm,columnsep=0.6cm]{geometry}
\usepackage{fancyhdr}
\usepackage{titlesec}
\usepackage{booktabs}
\usepackage{amsmath,amssymb,amsfonts}
\usepackage{microtype}
\usepackage{xcolor}
\usepackage{mdframed}
\usepackage{parskip}
\usepackage{enumitem}
\usepackage{hyperref}

\definecolor{rulered}{RGB}{180,30,30}
\definecolor{boxgray}{RGB}{245,245,245}
\definecolor{keywordgray}{RGB}{100,100,100}

\pagestyle{fancy}
\fancyhf{}
\fancyhead[L]{\textsc{\small Claude Mag}}
\fancyhead[C]{\small\textit{Machine Learning Research Digest}}
\fancyhead[R]{\small 2026-09-16}
\fancyfoot[C]{\small\thepage}
\renewcommand{\headrulewidth}{0.8pt}

\titleformat{\section}{\normalsize\bfseries\color{rulered}}{}{0pt}{}%
  [\vspace{-4pt}\color{rulered}\rule{\columnwidth}{0.4pt}]
\titlespacing{\section}{0pt}{8pt}{4pt}

\hypersetup{colorlinks=true,urlcolor=rulered,linkcolor=black}

\begin{document}
\setlength{\parindent}{0pt}
\setlength{\parskip}{4pt}

{\small\color{keywordgray}\textbf{Keywords:}
  heterogeneous hardware \textbullet{}
  subquadratic attention \textbullet{}
  disaggregation \textbullet{}
  energy efficiency \textbullet{}
  LLM serving}\par\vspace{4pt}

{\LARGE\bfseries\raggedright
  Rethinking Heterogeneous System
  Disaggregation for Subquadratic Attention}\par\vspace{4pt}

{\large\itshape\raggedright
  Prefill-Decode Disaggregation Misses Within-Decode Arithmetic
  Intensity Heterogeneity --- Fine-Grained SQD Delivers
  31--56\% Tokens-per-Joule Gains}\par\vspace{6pt}

{\small\textbf{By} Arya Tschand, Yaosheng Fu, Vikram Sharma Mailthody,
  Nicolai Oswald, Po-An Tsai, Ritchie Zhao, Oreste Villa,
  Vijay Janapa Reddi, Karu Sankaralingam
  (Harvard; NVIDIA; UW-Madison)}\par\vspace{2pt}

{\small\color{keywordgray}arXiv:2609.13134}
\par\vspace{2pt}\noindent{\color{rulered}\rule{\columnwidth}{0.6pt}}\vspace{6pt}

Frontier LLMs increasingly deploy subquadratic attention --- sparse,
linear, and sliding-window variants --- to reduce the $O(L)$ KV memory
footprint of long contexts. Heterogeneous hardware systems already
disaggregate by phase: prefill (compute-bound) on one device, decode
(memory-bound) on another. But this coarse split ignores a critical
observation: within a single decode step, different operators have
sharply different arithmetic intensity profiles. SQD (SubQuadratic
Disaggregation) exploits this by splitting the decode step itself,
routing memory-bound and compute-bound sub-phases to the hardware
optimized for each.

\begin{mdframed}[backgroundcolor=boxgray,linecolor=rulered,linewidth=1pt,
    innerleftmargin=6pt,innerrightmargin=6pt,
    innertopmargin=6pt,innerbottommargin=6pt]
\textbf{\small AT A GLANCE}\par\vspace{4pt}
\begin{description}[leftmargin=0pt,labelsep=4pt,itemsep=2pt,parsep=0pt,
    font=\small\bfseries]
  \item[Problem:] Prefill-decode disaggregation does not account for
    within-decode operator heterogeneity in models with subquadratic
    attention.
  \item[Why it matters:] 31--56\% energy efficiency gain is
    left on the table on frontier models (GLM 5.2, Nemotron 3 Ultra,
    Gemma 4 31B).
  \item[Method:] SQD splits decode by sub-phase: top-$k$ selection
    (memory device) vs.\ top-$k$ attention + FFN (compute device)
    for sparse; dense layers vs.\ subquadratic + FFN for linear models.
  \item[Result:] 53\%, 31\%, 56\% tokens/J improvement on GLM 5.2,
    Nemotron 3 Ultra, Gemma 4 31B on an 8$\times$B200 system.
  \item[Evidence:] Energy and throughput measured on production
    model weights; ablations by disaggregation level.
\end{description}
\end{mdframed}

\section{The Big Picture}

Full attention has arithmetic intensity $I = \text{FLOPs}/\text{bytes}
\approx d/2$ (roughly compute-bound for large $d$). Decode batch size 1
is memory-bound ($I \approx 1$) regardless of attention type. The
roofline model predicts runtime:
\[
  T(O) = \max\!\left(\frac{F}{C},\,\frac{B}{\beta}\right)
\]
where $C$ (FLOP/s) and $\beta$ (bytes/s) are device peaks. Routing
to the device where this bound is tightest minimizes runtime and
energy per token.

\section{Why Existing Approaches Fall Short}

Current disaggregation schemes split at the phase boundary: prefill
(high compute intensity, batch $= L$) goes to GPU; decode (low
intensity, batch $= 1$--few) goes to memory-optimized accelerator.
This is appropriate for dense attention. For subquadratic attention,
the decode step itself contains operators with radically different
intensities:

\begin{center}
\begin{tabular}{lcc}
\toprule
Operator & $I$ (FLOP/byte) & Bound \\
\midrule
Top-$k$ KV selection   & $\approx 1$     & Memory \\
Top-$k$ attn.\ + FFN   & $\approx 10^3$  & Compute \\
Dense attn.\ layer     & $\approx 10$    & Memory \\
Subquad.\ layer + FFN  & $\approx 2$     & Memory \\
\bottomrule
\end{tabular}
\end{center}

Routing the entire decode step to a memory-optimized device wastes
compute capacity on operators that are actually compute-bound.

\section{Core Mechanism: SQD}

\textbf{SQD for sparse attention.} Each decode step is split into two
sub-phases:
\begin{enumerate}[itemsep=2pt,topsep=2pt]
  \item \textit{Selection:} Traverse all $L$ keys to compute relevance
    scores and select top-$k$ entries. Cost: $O(L\cdot d_\text{head})$
    reads, $O(L)$ FLOPs. $I \approx 1$~FLOP/byte $\to$ memory device.
  \item \textit{Attention + FFN:} Attend to the selected $k$ entries
    and run the FFN. With $k \ll L$, KV footprint is small and static;
    $I \approx 10^3$~FLOP/byte $\to$ compute device.
\end{enumerate}

\textbf{SQD for linear/sliding-window attention.} Split by layer type:
\begin{enumerate}[itemsep=2pt,topsep=2pt]
  \item \textit{Dense attention layers:} Memory-bound at decode
    batch 1 $\to$ memory device.
  \item \textit{Subquadratic layers + FFN:} Linear attention state
    update is compute-bound $\to$ compute device.
\end{enumerate}

\textbf{Communication cost:} The two sub-phases exchange activation
tensors of size $O(d)$ per token per layer --- small compared with
$O(L\cdot d)$ KV data transfer.

\section{Technical Formulation}

The \textbf{ridge point} of a device is:
\[
  I^* = \frac{C}{\beta}
\]
On B200: $C = 4.5\times 10^{15}$~FP8~FLOP/s,
$\beta = 8\times 10^{12}$~bytes/s, so $I^* \approx 562$~FLOP/byte.

\textbf{Routing criterion:} Assign operator $O$ to the memory-optimized
device if $I(O) < I^*$, otherwise to the compute-optimized GPU.

\textbf{Energy per token:} Routing each operator to its optimal device
reduces average power draw per operation; the total energy per token is:
\[
  E = \sum_O \frac{F_O}{P_O} + \frac{B_O}{P_O^{\mathrm{mem}}}
\]
where $P_O$ is the TDP of the assigned device.

\section{Experimental Results}

Evaluated on an 8$\times$B200 heterogeneous system (4 compute-optimized
+ 4 memory-bandwidth-optimized) with production model weights:

\begin{center}
\begin{tabular}{llc}
\toprule
Model & Attn.\ Type & Tokens/J Gain \\
\midrule
GLM 5.2         & Sparse      & +53\% \\
Nemotron 3 Ultra & Linear     & +31\% \\
Gemma 4 31B     & Sliding-win & +56\% \\
\bottomrule
\end{tabular}
\end{center}

Throughput gains (tokens/s, batch $=$ 8, seq $=$ 32k):

\begin{center}
\begin{tabular}{lccc}
\toprule
Model & GPU-only & SQD & Speedup \\
\midrule
GLM 5.2     & 1,840 & 2,710 & 1.47$\times$ \\
Gemma 4 31B & 2,190 & 3,080 & 1.41$\times$ \\
\bottomrule
\end{tabular}
\end{center}

\section{Ablations and Interpretation}

\textbf{Selection-only disaggregation:} Routing only the top-$k$
selection phase accounts for 68\% of the total energy gain for sparse
attention models.

\textbf{Without within-decode layer-type split:} Routing by
prefill/decode phase only reduces energy gain from 53\% to 29\% for
GLM 5.2, quantifying the value of intra-decode disaggregation.

\textbf{Communication overhead:} Cross-device activation transfer
adds 7--12\% latency but is offset by 40--55\% reduction in
time-per-token on memory-bound operators.

\textbf{Dense models:} SQD provides no benefit when all attention
layers are dense --- standard prefill-decode disaggregation is
already optimal in that case.

\section*{Reference}

Arya Tschand et al.\ \textbf{Rethinking Heterogeneous System
Disaggregation for Subquadratic Attention.}
arXiv:2609.13134, September 2026.
\url{https://arxiv.org/abs/2609.13134}

\end{document}
